% 原创算子图；几何位置与前向模式相同，虚线反向传递敏感度。
\begin{tikzpicture}[x=1mm,y=1mm,font=\figurefont\footnotesize,text=NNInk,
  ad input/.style={circle,draw=NNInput,fill=NNInput!18,line width=1.1pt,
    minimum size=11mm,inner sep=0pt},
  ad operator/.style={circle,draw=NNHidden,fill=NNHidden!18,line width=1.1pt,
    minimum size=11mm,inner sep=0pt},
  ad reverse/.style={-{Stealth[length=1.8mm,width=1.3mm]},draw=NNGradient,
    line width=1.2pt,dashed,rounded corners=2mm},
  ad derivative/.style={text=NNGradient,align=center,text width=30mm,
    inner sep=0pt}]
  \node[ad input] (u) at (12,14) {$u=2$};
  \node[ad input] (v) at (12,-14) {$v=3$};
  \node[ad derivative] at (12,3) {$\bar u=32$};
  \node[ad derivative] at (12,-25) {$\bar v=16$};
  \node[ad operator] (mul) at (52,0) {$\times$};
  \node[ad operator] (add) at (84,0) {$+$};
  \node[ad operator] (square) at (116,0) {$(\cdot)^2$};
  \node[ad operator,draw=NNOutput,fill=NNOutput!18] (half) at (148,0)
    {$\times\frac12$};
  \draw[ad reverse] (mul.north west)--
    node[above,sloped,text=NNGradient] {$8\cdot3=24$} (u.east);
  \draw[ad reverse] (mul.south west)--
    node[below,sloped,text=NNGradient] {$8\cdot2=16$} (v.east);
  \draw[ad reverse] (add.north)--(84,28)--(12,28)--(u.north);
  \node[above,text=NNGradient] at (49,28) {$8\cdot1=8$};
  \draw[ad reverse] (add)--node[above,text=NNGradient] {$\times1$} (mul);
  \draw[ad reverse] (square)--node[above,text=NNGradient] {$\times2s$} (add);
  \draw[ad reverse] (half)--node[above,text=NNGradient] {$\times\tfrac12$} (square);
  \foreach \x/\value in {52/r=6,84/s=8,116/t=64,148/J=32}{
    \node[text=NNInput] at (\x,-12) {$\value$};
    \draw[NNLine,line width=.45pt] (\x-12,-20)--(\x+12,-20);
  }
  \node[ad derivative] at (52,-30) {$\bar r=8$};
  \node[ad derivative] at (84,-30) {$\bar s=8$};
  \node[ad derivative] at (116,-30) {$\bar t=\tfrac12$};
  \node[ad derivative] at (148,-30) {$\bar J=1$\\求导起点};
  \draw[ad reverse] (32,-47)--(22,-47);
  \node[anchor=west] at (35,-47) {梯度回传};
  \node[anchor=west,text=NNGradient] at (84,-47) {$\bar u=24+8=32$};
  \node[anchor=west] at (125,-47) {分支贡献相加};
\end{tikzpicture}
